The European Copernicus Earth-observation programme operates two Sentinel-2
remote sensing satellites. These satellites orbit the Earth on Sun-synchronous
polar orbits at about 800\,km altitude. They pass over a given area once every
few days, always taking images at the same local time (accurate to within a few
minutes). The cameras are sensitive to 13 different optical and near-infrared
spectral bands. The resolution of the images is 10 meters.

The 3rd tallest building in Hungary is the chimney of a power plant near the
town Tisza\'ujv\'aros. You can see two Sentinel-2 satellite images. One is from
June 29, another one is from December 16, close to the summer and winter
solstices, respectively. The orientation of the images is as normal, i.e.\
north is up and east is to the right.

\ioaafig{2019_TH_6-f1.pdf}

The estimated shadow lengths, based on the images above and the scales given in
the lower-left corners, of them are $x_1 = 125$\,m and $x_2 = 780$\,m. Answer
the following questions:

\begin{description}
\item[a)] On which date do we expect the shadow to be longer? \textbf{Mark the
  letter of your answer with $\times$ on the answer sheet.} \hfill[1\,p]\\
  (A) on June 29\\
  (B) on December 16
\item[b)] At which time of the day did the Sentinel-2 satellites fly over this
  area? \textbf{Mark the letter of your answer with $\times$ on the answer
  sheet.} \hfill[1\,p]\\
  (A) early morning\\
  (B) late morning\\
  (C) early afternoon\\
  (D) late afternoon
\item[c)] Based on the given shadow lengths, estimate the height of this
  chimney. For this calculation only, assume that the satellite images were
  taken at local noon. \hfill[16\,p]
\item[d)] What could affect the accuracy of the chimney's derived height (more
  than one choice is possible)? \textbf{Mark the letter of your answer with
  $\times$ on the answer sheet.} \hfill[2\,p]\\
  (A) The oblate spheroid shape of the Earth.\\
  (B) Limited resolution of the satellite images and the ill-defined edge of
  the shadow.\\
  (C) The elevation of the base of the tower above sea level.\\
  (D) Seasonal variation in the tilt of the Earth's rotational axis.\\
  (E) Taking into account the effect of atmospheric refraction.
\end{description}
